\documentclass[11pt]{article}

\usepackage[margin=1.0in]{geometry}
\usepackage{lecnotes}
\input{lmacros}

\newcommand{\course}{15-316: Software Foundations of Security \& Privacy}
\newcommand{\lecturer}{}
\newcommand{\lecdate}{Due Tue, Mar 24, 2026}
\newcommand{\lecnum}{5}
\newcommand{\lectitle}{Homework 5}
\newcommand{\courseurl}{https://15316-cmu.github.io/2026/}

% HW1 : 70 pts
% HW2 : 80 pts
% HW3 : 75 pts
% HW4 : 90 pts + 20 pts EC
% HW5 : 85 pts
% ------------
%      400 pts


\title{Assignment \lecnum \\ Proof-Carrying Authorization}
\date{\lecdate \\ 85 points}

\lhead[\fancyplain{}{\bfseries A\lecnum.\thepage}]%
      {\fancyplain{}{\bfseries\lectitle}}
\chead[]{}
\rhead[\fancyplain{}{\bfseries\lectitle}]%
      {\fancyplain{}{\bfseries A\lecnum.\thepage}}
\lfoot[{\small\scshape Assignments}]{{\small\scshape Assignments}}
\cfoot[]{}
\rfoot[{\small\scshape\lecdate}]{{\small\scshape\lecdate}}

\begin{document}

\maketitle

Your solution should be handed in as a file \texttt{hw5.pdf} to Gradescope.  If
at all possible, write your solutions in \LaTeX.  The handout
\texttt{hw5-pca.zip} includes the \LaTeX\ sources for some lectures and the
necessary style files which provide some examples for rules, derivations, and
proofs.  You may use up to two late days as usual.

\section{Authorization Logic [25 points]}

For each of the following, either give a derivation in the sequent calculus or
give a counterexample.  In the latter case, use $p$ and $q$ as atomic formulas
and demonstrate that the counterexample is not derivable.

\begin{task}[5 pts]
 $(A \says P) \land (A \says Q) \vdash A \says (P \land Q)$
\end{task}

\begin{task}[5 pts]
 $A \says (P \land Q) \vdash (A \says P) \land (A \says Q)$
\end{task}

\begin{task}[5 pts]
 $(A \says P) \lor (A \says Q) \vdash A \says (P \lor Q)$
\end{task}

\begin{task}[5 pts]
 $A \says (P \lor Q) \vdash (A \says P) \lor (A \says Q)$
\end{task}

\begin{task}[5 pts]
 $A \says \forall x.\, P(x) \vdash \forall x.\, A \says P(x)$
\end{task}

\section{Trust [30 points]}
\newcommand{\trusts}{\mathrel{\mb{trusts}}}

We informally define that $A \trusts B$ if whenever $B$ affirms $P$ then $A$
also affirms $P$.  Unfortunately, we cannot express this directly in our
authorization logic because the right-hand side of the definition
\[
  A \trusts B \triangleq \forall P.\, B \says P \arrow A \says P
\]
quantifies over all possible propositions.  In authorization logic, quantifiers
only range over principals and constants from some finite domain (like rooms or
files or homework assignments, etc.).

Instead we assume that there is a \emph{fixed preorder of principals} where
$A \leq B$ means that $A$ trusts $B$.  We assume that this relation is reflexive
and transitive.

\begin{task}[15 pts]\rm
  Define rules in the sequent calculus for authorization logic that incorporate
  the trust relationship.  Your premises may directly refer to $A \leq B$ for
  principals $A$ and $B$.  You may modify existing rules ${\says}R$, ${\says}L$,
  and ${\aff}$, and you may add new rules.  Other inference rules should
  remain unchanged.
\end{task}

\begin{task}[5 pts]\rm
  Prove $A \says P \land B \says Q \vdash A \says (P \land Q)$ in your system,
  provided $A \leq B$.
\end{task}

\begin{task}[10 pts]\rm
  Give a counterexample to
  $A \says P \land B \says Q \vdash B \says (P \land Q)$ provided
  $B \not\leq A$.  You should be able to use your rules to prove that the
  counterexample cannot be derived.  If not, explain for partial credit why it
  is difficult or impossible.
\end{task}

\section{Groups and Access Control [30 points]}

In the AFS file system used on the \verb'linux.andrew' machines there are two
kinds of principals: individuals (identified by their Andrew id) and groups
(identified by a name distinct from an individual id).  We formalize a variant
of the Andrew file system access control with the following policies:
\begin{enumerate}[(i)]
\item Any principal may create a new group they own simply by saying so.  Assume
  that adding such an affirmation to the policy would fail if the group already
  exists.
\item The owner of a group is always one of its members.
\item Only the owner of a group may add members to it.  We are not concerned
  with revoking group membership.
\item The owner of a file may grant a principal read or write access simply by
  saying so.
\item If a group has read or write access to a file, each member of the group
  has read or write access, respectively, to the file.
\end{enumerate}
In order to formalize this policy in authorization logic, we
use the following vocabulary:
\begin{itemize}
\item Principals $A, B, C, \ldots$ could denote an individual or a group.
\item Groups $G, \ldots$ if we know a principal must be a group.
\item $\mi{admin}$. The principal who administers policy, affirming rules for
  groups and access control.  Depending on the configuration, $\mi{admin}$ could be a
  group or an individual.
\item Mode $M$, which is either $\mi{read}$ or $\mi{write}$.
\item File names $F$.
\item $\ms{ownsGroup}(A, G)$.  Principal $A$ owns group $G$.
\item $\ms{ownsFile}(A, F)$.  Principal $A$ owns file $F$.
\item $\ms{member}(A, G)$.  Principal $A$ is a member of group $G$.
\item $\ms{mayOpen}(A, F, M)$.  Principal $A$ may open file $F$ in mode $M$.
\end{itemize}

\begin{task}[10 points]\rm
  Write out the access control policy explained above in authorization logic.
  Label your affirmations so they can be used in proof terms.
\end{task}

\begin{task}[10 points]\rm
  Write out the necessary affirmations by $\mi{mfredrik}$ or $\mi{admin}$ to
  express each of the following.  Label your affirmations so they
  can be used in proof terms.
  \begin{enumerate}
  \item Individual $\mi{mfredrik}$ owns group $\mi{316\_ta}$
    and files $\mi{316\_roster}$ and $\mi{316\_grades}$.
  \item Individuals $\mi{kmann}$ and $\mi{jzou}$ are members of $\mi{316\_ta}$.
  \item Group $\mi{316\_ta}$ may read file $\mi{316\_roster}$ and read and write file
    $\mi{316\_grades}$.
  \end{enumerate}
\end{task}

\begin{task}[10 points]\rm
  Write out a proof term (as defined in
  \href{\courseurl/lectures/18-proofrep.pdf}{Lecture 18}) showing $\mi{kmann}$
  may write file $\mi{316\_grades}$.  You should use labels from the
  affirmations in the two previous tasks.
\end{task}

\end{document}
